\section{Checkpoint Evaluation}
\label{app:pending-evaluations}

The common inference benchmark uses the same GPU, precision, and
measurement procedure for all methods. Each method retains its
accuracy-evaluation attention setting, and achieved token counts
are recorded. The benchmark reports
full-validation accuracy separately
from batch-1 latency, batch-64 throughput, and peak allocated memory.
Table~\ref{tab:systems-pending} gives the trained-checkpoint measurements;
the synthetic random-weight timings in
Appendices~\ref{app:synthetic} and \ref{app:followup} are
reported separately.

\begin{table}[h]
\caption{Common-protocol trained-checkpoint inference at 224 pixels.
B1/B64 denote batch sizes 1 and 64.}
\label{tab:systems-pending}
\centering\footnotesize
\setlength{\tabcolsep}{2.8pt}
\begin{tabular}{lrrrr}
\toprule
Method & Patches & B1 (ms) & B64 (im/s) & B64 alloc. (GiB) \\
\midrule
Dense DeiT-S/8 & 784 & \pendingvalue & \pendingvalue & \pendingvalue \\
DTEM-p8 & 392 & \pendingvalue & \pendingvalue & \pendingvalue \\
DeiT-S/8 + ToMe & 392 & \pendingvalue & \pendingvalue & \pendingvalue \\
DeiT-S/8 + PiToMe & 392 & \pendingvalue & \pendingvalue & \pendingvalue \\
\mn{} ($R=3$) & 392 & \pendingvalue & \pendingvalue & \pendingvalue \\
\bottomrule
\end{tabular}

\end{table}

The component test uses the same trained \mn{} EMA checkpoint
for all four rows of Table~\ref{tab:components}. The mass-bias
intervention sets only the latent log-mass key bias to zero;
the recovery intervention bypasses only residual cross-attention.
The native score is checked against the common inference evaluation.
Paired intervals compare per-image correctness on the same
50,000 validation examples and do not represent training-seed
uncertainty.

\begin{table}[h]
\caption{Inference-only component test at 224 pixels.
$\Delta$ and its paired 95\% interval are relative to native top-1.}
\label{tab:components}
\centering\small
\begin{tabular}{lrr}
\toprule
Inference setting & Top-1 (\%) & $\Delta$ top-1 (pp), 95\% CI \\
\midrule
Native (mass + recovery) & \pendingvalue & --- \\
Mass bias off & \pendingvalue & \pendingvalue \\
Recovery off & \pendingvalue & \pendingvalue \\
Both off & \pendingvalue & \pendingvalue \\
\bottomrule
\end{tabular}

\end{table}
